\documentclass[11pt, a4paper]{article}

\usepackage[a4paper, margin=2.5cm]{geometry}
\usepackage{fontspec}
\usepackage[russian, bidi=basic, provide=*]{babel}
\babelprovide[import, onchar=ids fonts]{russian}
\babelprovide[import, onchar=ids fonts]{english}
\babelfont{rm}{Noto Sans}
\babelfont{sf}{Noto Sans}
\babelfont{tt}{Noto Sans Mono}

\usepackage{enumitem}
\setlist[itemize]{label=-}
\usepackage{amsmath}
\usepackage{booktabs}
\usepackage{tabularx}
\usepackage{listings}
\usepackage{xcolor}
\usepackage[htt]{hyphenat}
\usepackage{microtype}
\usepackage{hyperref}

\setlength{\emergencystretch}{3em}

\definecolor{codegray}{rgb}{0.5,0.5,0.5}
\definecolor{backcolor}{rgb}{0.97,0.97,0.97}

\lstset{
    backgroundcolor=\color{backcolor},
    basicstyle=\ttfamily\scriptsize,
    breaklines=true,
    breakatwhitespace=true,
    breakautoindent=true,
    postbreak=\mbox{\textcolor{codegray}{\(\hookrightarrow\)\ }},
    columns=fullflexible,
    keepspaces=true,
    frame=single,
    rulecolor=\color{black!20},
    numbers=left,
    numbersep=6pt,
    numberstyle=\tiny\color{codegray},
    showstringspaces=false,
    tabsize=2,
    xleftmargin=15pt,
    framexleftmargin=15pt
}

\hypersetup{
    colorlinks=true,
    linkcolor=black,
    urlcolor=blue
}

\begin{document}

\begin{titlepage}
    \centering
    {\small МИНИСТЕРСТВО НАУКИ И ВЫСШЕГО ОБРАЗОВАНИЯ РОССИЙСКОЙ ФЕДЕРАЦИИ\par}
    \vspace{0.3cm}
    {\small Федеральное государственное автономное образовательное учреждение высшего образования\par}
    {\small «Московский авиационный институт (национальный исследовательский университет)»\par}
    
    \vspace{4cm}
    {\Large \textbf{ОТЧЁТ ПО ЛАБОРАТОРНОЙ РАБОТЕ №4}\par}
    
    \vspace{1.2cm}
    {\large \textbf{Дисциплина:}\par}
    {\large Операционные системы и архитектура компьютера\par}
    
    \vspace{0.8cm}
    {\large \textbf{Тема:}\par}
    {\large Освоение принципов работы с файловыми системами. Обеспечение обмена данными между процессами посредством технологии «File mapping»\par}
    
    \vspace{3.5cm}
    \begin{flushright}
        \begin{tabular}{@{}r@{}}
            \textbf{Выполнил:} \\
            Студент группы М8О-203БВ-25 \\
            К. П. Смаге \\[0.8cm]
            \textbf{Преподаватель:} \\
            Е. С. Миронов \\
        \end{tabular}
    \end{flushright}
    
    \vfill
    {\large Москва, 2026 г.\par}
\end{titlepage}

\section{Цель работы}
Приобретение практических навыков в:
\begin{itemize}
    \item освоении принципов работы с файловыми системами и отображением объектов в виртуальную память процесса в ОС Linux (POSIX);
    \item обеспечении межпроцессного взаимодействия (IPC) посредством технологии «File mapping» (разделяемой памяти \texttt{shm\_open} + \texttt{mmap}) и синхронизации доступа с использованием POSIX-семафоров.
\end{itemize}

\section{Краткое описание архитектуры решения}
Разработано решение работы для ОС семейства Linux (POSIX). Программный комплекс состоит из двух отдельных исполняемых файлов: управляющей родительской программы (\texttt{parent}) и рабочей дочерней программы (\texttt{child}).

Алгоритм функционирования системы включает следующие этапы:
\begin{enumerate}
    \item Родительский процесс устанавливает обработчики сигналов \texttt{SIGINT} и \texttt{SIGTERM} (через \texttt{sigaction} без флага \texttt{SA\_RESTART}) для гарантированной очистки ресурсов ОС при прерывании, после чего считывает со стандартного потока ввода два различных пути к выходным файлам для первого и второго дочерних процессов.
    \item Родительский процесс инициализирует два экземпляра класса \texttt{Process} (\texttt{child1} и \texttt{child2}). Для каждого дочернего процесса до вызова \texttt{fork()} создается отдельный объект разделяемой памяти с уникальным именем (содержащим PID родителя и порядковый индекс) через \texttt{shm\_open} с флагами \texttt{O\_CREAT | O\_EXCL | O\_RDWR} и правами \texttt{0600}, задается его размер через \texttt{ftruncate()}, выполняется отображение в адресное пространство с флагом \texttt{MAP\_SHARED} через \texttt{mmap()}, после чего файловый дескриптор закрывается. Внутри отображенной области инициализируются два неименованных межпроцессных семафора (\texttt{ready} и \texttt{ack}, \texttt{pshared = 1}).
    \item Родительский процесс открывает целевой выходной файл на запись (\texttt{open()}), порождает дочерний процесс (\texttt{fork()}), перенаправляет его стандартный поток вывода (\texttt{stdout}) в открытый файл через \texttt{dup2()} и запускает программу \texttt{child} через \texttt{execl()}, передавая имя объекта разделяемой памяти аргументом командной строки. Стандартный поток ошибок (\texttt{stderr}) остается привязанным к терминалу для вывода диагностических сообщений и идентификаторов процессов (\texttt{PID}).
    \item Родительский процесс построчно считывает пользовательский ввод произвольной длины и осуществляет фильтрацию согласно условию варианта 20:
    \begin{itemize}
        \item если длина строки превышает 10~символов ($\text{length} > 10$), строка пересылается процессу \texttt{child2} через его сегмент разделяемой памяти;
        \item если длина строки не превышает 10~символов ($\text{length} \le 10$), строка пересылается процессу \texttt{child1} через его сегмент разделяемой памяти.
    \end{itemize}
    Для поддержки строк произвольной длины без выхода за границы отображенной области передача организована порциями до \texttt{CHUNK\_SIZE} (4096~байт) с заголовком длины и флагом пакета (\texttt{DATA\_CHUNK}, \texttt{LINE\_END}, \texttt{END\_OF\_FILE}). Синхронизация каждого пакета выполняется по схеме подтверждения: запись данных $\rightarrow$ \texttt{sem\_post(\&ready)} $\rightarrow$ ожидание подтверждения \texttt{sem\_timedwait(\&ack)} с проверкой жизнеспособности потомка через \texttt{waitid(..., WNOHANG | WNOWAIT)}.
    \item Каждый дочерний процесс открывает готовый объект разделяемой памяти (\texttt{shm\_open} с \texttt{O\_RDWR}), проверяет его размер через \texttt{fstat()} (для предотвращения \texttt{SIGBUS}) и отображает область через \texttt{mmap(..., MAP\_SHARED, ...)}. В цикле потомок ожидает \texttt{sem\_wait(\&ready)}, проверяет корректность длины чанка и флага, накапливает строку, подтверждает чтение через \texttt{sem\_post(\&ack)}, а по получении флага \texttt{LINE\_END} инвертирует строку и записывает результат в перенаправленный \texttt{stdout}.
    \item При получении конца ввода (\texttt{EOF}) или сигнала завершения родительский процесс отправляет обоим потомкам пакет с флагом \texttt{END\_OF\_FILE}, ожидает их завершения через \texttt{waitpid()} (исключая образование процессов-зомби), после чего уничтожает семафоры (\texttt{sem\_destroy}), снимает отображение (\texttt{munmap}) и удаляет имена объектов из \path{/dev/shm} (\texttt{shm\_unlink}). Дочерние процессы при завершении выполняют только снятие собственного отображения (\texttt{munmap}).
\end{enumerate}

Поскольку виртуальные адреса отображения в разных процессах не совпадают, структура \texttt{SharedRegion} спроектирована плоской и не содержит указателей или динамических контейнеров — в общей памяти хранятся только семафоры, длина порции, флаг состояния и фиксированный буфер байтов. Логика приложения модульно разделена на транспортный слой (\path{include/shm_channel.hpp}, \path{src/shm_channel.cpp}), управление жизненным циклом процессов (\path{include/process.hpp}, \path{src/process.cpp}), утилиты обработки строк (\path{include/string_utils.hpp}, \path{src/string_utils.cpp}) и точки входа (\path{src/parent.cpp}, \path{src/child.cpp}). Оба системных класса реализуют идиому RAII и проверяют коды возврата всех системных вызовов с выводом диагностики через \texttt{perror()}.

В таблице 1 приведено сопоставление используемых системных вызовов и примитивов POSIX с их ролью в жизненном цикле программы.

\begin{table}[h!]
\centering
\small
\begin{tabularx}{\textwidth}{l >{\raggedright\arraybackslash}X >{\raggedright\arraybackslash}X}
\toprule
\textbf{Операция ОС} & \textbf{Примитив POSIX} & \textbf{Назначение в решении} \\
\midrule
Создание shm & \texttt{shm\_open(name, flags, mode)} & Создание уникального объекта (\texttt{O\_CREAT | O\_EXCL | O\_RDWR}, \texttt{0600}) в родителе и подключение (\texttt{O\_RDWR}) в потомке \\
Размер сегмента & \texttt{ftruncate(fd, size)}, \texttt{fstat(fd, \&st)} & Установка размера сегмента в родителе и контроль границ перед отображением в потомке \\
Отображение & \texttt{mmap(..., PROT\_READ | PROT\_WRITE, MAP\_SHARED, fd, 0)} & Связывание общих физических страниц памяти с виртуальными адресными пространствами процессов \\
Синхронизация & \texttt{sem\_init}, \texttt{sem\_post}, \texttt{sem\_wait}, \texttt{sem\_timedwait} & Упорядочивание записи и чтения (\texttt{ready}/\texttt{ack}) без активного ожидания и защита от зависания при падении потомка \\
Файлы и потоки & \texttt{open()}, \texttt{dup2(fd, STDOUT\_FILENO)}, \texttt{close()} & Открытие целевых выходных файлов и привязка \texttt{stdout} дочерних процессов к файлам \\
Процессы & \texttt{fork()}, \texttt{execl()}, \texttt{waitid()}, \texttt{waitpid()} & Порождение потомков, передача имени сегмента, контроль состояния и сбор статусов завершения без зомби-процессов \\
Сигналы & \texttt{sigaction(SIGINT / SIGTERM, ...)} & Корректное прерывание ввода с гарантированным освобождением сегментов разделяемой памяти \\
Очистка ресурсов & \texttt{sem\_destroy()}, \texttt{munmap()}, \texttt{shm\_unlink()} & Раздельное снятие отображения в каждом процессе и удаление имен объектов из \path{/dev/shm} в родителе \\
\bottomrule
\end{tabularx}
\caption{Используемые системные вызовы и примитивы POSIX API}
\end{table}

\section{Листинг ключевых фрагментов кода}

\subsection{Структура разделяемой памяти и интерфейс канала IPC}
Заголовочный файл \path{include/shm_channel.hpp} определяет плоскую структуру \texttt{SharedRegion} (без указателей), формат пакетов для порционной передачи строк произвольной длины и RAII-класс \texttt{ShmChannel}.

\begin{lstlisting}[language=C++]
constexpr std::size_t CHUNK_SIZE = 4096;

enum class PacketFlag : std::uint8_t {
    DATA_CHUNK = 0,
    LINE_END = 1,
    END_OF_FILE = 2
};

struct SharedRegion {
    sem_t ready;
    sem_t ack;
    std::size_t length;
    PacketFlag flag;
    char data[CHUNK_SIZE];
};

class ShmChannel {
public:
    ShmChannel() = default;
    ~ShmChannel();

    ShmChannel(const ShmChannel&) = delete;
    ShmChannel& operator=(const ShmChannel&) = delete;

    ShmChannel(ShmChannel&& other) noexcept;
    ShmChannel& operator=(ShmChannel&& other) noexcept;

    bool create(const std::string& name);
    bool open(const std::string& name);

    bool send_line(const std::string& line, pid_t peer_pid = -1);
    bool send_eof(pid_t peer_pid = -1);
    bool receive_line(std::string& out_line, bool& is_eof);

    void close();

    bool is_open() const noexcept;
    const std::string& get_name() const noexcept;

private:
    std::string shm_name_;
    SharedRegion* region_ = nullptr;
    bool is_owner_ = false;
    bool sems_initialized_ = false;
};
\end{lstlisting}
\begin{center}
    \small Листинг 1 — Структура разделяемой памяти и объявление класса \texttt{ShmChannel} (\path{include/shm_channel.hpp})
\end{center}

\subsection{Инициализация разделяемой памяти и синхронизация обмена}
В файле \path{src/shm_channel.cpp} реализованы создание объекта разделяемой памяти, контроль границ при подключении, порционная отправка с ожиданием подтверждения (\texttt{sem\_timedwait} + \texttt{waitid} с \texttt{WNOWAIT}) и прием строк с проверкой длины и флагов.

\begin{lstlisting}[language=C++]
bool ShmChannel::create(const std::string& name) {
    if (is_open()) {
        close();
    }

    int fd = shm_open(name.c_str(), O_CREAT | O_EXCL | O_RDWR, 0600);
    if (fd == -1) {
        perror("shm_open");
        return false;
    }

    if (ftruncate(fd, static_cast<off_t>(sizeof(SharedRegion))) == -1) {
        perror("ftruncate");
        ::close(fd);
        shm_unlink(name.c_str());
        return false;
    }

    void* addr = mmap(
        nullptr, sizeof(SharedRegion),
        PROT_READ | PROT_WRITE, MAP_SHARED, fd, 0
    );
    if (addr == MAP_FAILED) {
        perror("mmap");
        ::close(fd);
        shm_unlink(name.c_str());
        return false;
    }

    if (::close(fd) == -1) {
        perror("close");
        munmap(addr, sizeof(SharedRegion));
        shm_unlink(name.c_str());
        return false;
    }

    shm_name_ = name;
    region_ = static_cast<SharedRegion*>(addr);
    is_owner_ = true;
    sems_initialized_ = false;

    region_->length = 0;
    region_->flag = PacketFlag::DATA_CHUNK;
    std::memset(region_->data, 0, sizeof(region_->data));

    if (sem_init(&region_->ready, 1, 0) == -1) {
        perror("sem_init ready");
        close();
        return false;
    }
    if (sem_init(&region_->ack, 1, 0) == -1) {
        perror("sem_init ack");
        sem_destroy(&region_->ready);
        close();
        return false;
    }

    sems_initialized_ = true;
    return true;
}

bool ShmChannel::send_line(const std::string& line, pid_t peer_pid) {
    if (!is_open()) {
        return false;
    }

    if (line.empty()) {
        region_->length = 0;
        region_->flag = PacketFlag::LINE_END;
        if (sem_post(&region_->ready) == -1) {
            perror("sem_post ready");
            return false;
        }
        return wait_ack(&region_->ack, peer_pid);
    }

    std::size_t offset = 0;
    while (offset < line.size()) {
        std::size_t remaining = line.size() - offset;
        std::size_t chunk_len = (remaining > CHUNK_SIZE) ? CHUNK_SIZE : remaining;

        std::memcpy(region_->data, line.data() + offset, chunk_len);
        region_->length = chunk_len;
        offset += chunk_len;
        region_->flag = (offset >= line.size())
            ? PacketFlag::LINE_END
            : PacketFlag::DATA_CHUNK;

        if (sem_post(&region_->ready) == -1) {
            perror("sem_post ready");
            return false;
        }
        if (!wait_ack(&region_->ack, peer_pid)) {
            return false;
        }
    }

    return true;
}

bool ShmChannel::receive_line(std::string& out_line, bool& is_eof) {
    out_line.clear();
    is_eof = false;

    if (!is_open()) {
        return false;
    }

    while (true) {
        if (!wait_sem(&region_->ready)) {
            return false;
        }

        std::size_t len = region_->length;
        PacketFlag flag = region_->flag;

        if (len > CHUNK_SIZE) {
            sem_post(&region_->ack);
            return false;
        }
        if (flag == PacketFlag::END_OF_FILE) {
            is_eof = true;
            sem_post(&region_->ack);
            return false;
        }
        if (flag != PacketFlag::DATA_CHUNK && flag != PacketFlag::LINE_END) {
            sem_post(&region_->ack);
            return false;
        }

        if (len > 0) {
            out_line.append(region_->data, len);
        }
        if (sem_post(&region_->ack) == -1) {
            perror("sem_post ack");
            return false;
        }
        if (flag == PacketFlag::LINE_END) {
            return true;
        }
    }
}
\end{lstlisting}
\begin{center}
    \small Листинг 2 — Создание общей памяти и синхронизированная передача пакетов (\path{src/shm_channel.cpp})
\end{center}

\subsection{Порождение дочернего процесса и управление ресурсами}
В файле \path{src/process.cpp} объект разделяемой памяти инициализируется до вызова \texttt{fork()}, что исключает гонку на старте. В дочерней ветке при ошибке используется \texttt{\_exit()}, чтобы не вызывать деструктор родительского канала, а в методе \texttt{wait()} сегмент освобождается строго после завершения потомка.

\begin{lstlisting}[language=C++]
Process::Process(
    const std::string& binary_path,
    const std::string& output_file_path
)
    : pid_(-1), is_closed_(false) {
    unsigned long index = next_shm_index.fetch_add(1, std::memory_order_relaxed);
    std::string shm_name =
        "/ipc_shm_" + std::to_string(getpid()) + "_" + std::to_string(index);

    if (!channel_.create(shm_name)) {
        throw std::runtime_error("Failed to create shared memory channel: " + shm_name);
    }

    int file_fd = open(
        output_file_path.c_str(),
        O_WRONLY | O_CREAT | O_TRUNC,
        0644
    );
    if (file_fd == -1) {
        perror("open");
        channel_.close();
        throw std::runtime_error("Failed to open output file: " + output_file_path);
    }

    std::cout.flush();
    std::cerr.flush();

    pid_t fork_rc = fork();
    if (fork_rc == -1) {
        perror("fork");
        close(file_fd);
        channel_.close();
        throw std::runtime_error("Failed to fork");
    } else if (fork_rc == 0) {
        if (file_fd != STDOUT_FILENO) {
            if (dup2(file_fd, STDOUT_FILENO) == -1) {
                perror("dup2 stdout");
                _exit(EXIT_FAILURE);
            }
            if (close(file_fd) == -1) {
                perror("close file_fd");
                _exit(EXIT_FAILURE);
            }
        }

        execl(
            binary_path.c_str(),
            binary_path.c_str(),
            shm_name.c_str(),
            static_cast<char*>(nullptr)
        );

        perror("execl");
        _exit(127);
    } else {
        pid_ = fork_rc;
        if (close(file_fd) == -1) {
            perror("close file_fd");
        }
    }
}

void Process::wait() {
    if (pid_ > 0) {
        if (!is_closed_) {
            close_input();
        }

        int status = 0;
        while (waitpid(pid_, &status, 0) == -1) {
            if (errno == EINTR) {
                continue;
            }
            perror("waitpid");
            break;
        }
        pid_ = -1;
    }

    channel_.close();
}
\end{lstlisting}
\begin{center}
    \small Листинг 3 — Инициализация IPC, порождение процесса и ожидание завершения (\path{src/process.cpp})
\end{center}

\subsection{Маршрутизация данных и обработка сигналов в родительском процессе}
Функция \texttt{parent()} в файле \path{src/parent.cpp} настраивает обработчики сигналов \texttt{SIGINT}/\texttt{SIGTERM}, считывает имена двух выходных файлов, запускает оба дочерних процесса, распределяет входные строки по критерию длины ($\le 10$ или $> 10$~символов) и завершает работу без утечек ресурсов.

\begin{lstlisting}[language=C++]
void setup_signal_handlers() {
    struct sigaction sa {};
    sa.sa_handler = handle_signal;
    sigemptyset(&sa.sa_mask);
    sa.sa_flags = 0;

    if (sigaction(SIGINT, &sa, nullptr) == -1) {
        perror("sigaction SIGINT");
        throw std::runtime_error("Failed to install SIGINT handler");
    }
    if (sigaction(SIGTERM, &sa, nullptr) == -1) {
        perror("sigaction SIGTERM");
        throw std::runtime_error("Failed to install SIGTERM handler");
    }
}

void parent() {
    setup_signal_handlers();

    std::cerr << "[Parent] PID: " << get_current_pid() << '\n';

    std::cout << "Enter output file for Child1: ";
    std::string file1{};
    if (!std::getline(std::cin, file1) || file1.empty()) {
        if (g_stop_requested) {
            return;
        }
        throw std::runtime_error("Failed to read filename for Child1");
    }

    std::cout << "Enter output file for Child2: ";
    std::string file2{};
    if (!std::getline(std::cin, file2) || file2.empty()) {
        if (g_stop_requested) {
            return;
        }
        throw std::runtime_error("Failed to read filename for Child2");
    }

    if (file1 == file2) {
        throw std::runtime_error("Output files for Child1 and Child2 must be different");
    }

    Process child1(CHILD_PATH, file1);
    Process child2(CHILD_PATH, file2);

    std::cerr << "[Parent] Created Child1 (PID: " << child1.get_pid() << ")\n";
    std::cerr << "[Parent] Created Child2 (PID: " << child2.get_pid() << ")\n";

    std::string line{};
    std::cout << "Enter lines for processing (Ctrl+D / Ctrl+C to finish): \n";
    while (!g_stop_requested && std::getline(std::cin, line)) {
        if (line.length() > 10) {
            child2.send_line(line);
        } else {
            child1.send_line(line);
        }
    }

    child1.close_input();
    child2.close_input();

    child1.wait();
    child2.wait();
}
\end{lstlisting}
\begin{center}
    \small Листинг 4 — Логика диспетчеризации строк и обработки сигналов (\path{src/parent.cpp})
\end{center}

\subsection{Прием и обработка строк в дочернем процессе}
Дочерний процесс (\path{src/child.cpp}) подключается к разделяемой памяти по имени, полученному через \texttt{argv[1]}, принимает строки через \texttt{ShmChannel::receive\_line()}, инвертирует их функцией \texttt{reverse\_string()} и записывает результат в перенаправленный поток \texttt{std::cout} с контролем ошибок вывода.

\begin{lstlisting}[language=C++]
void reverse_string(std::string& line) {
    std::reverse(line.begin(), line.end());
}

int child(const std::string& shm_name) {
    std::cerr << "[Child] PID: " << static_cast<unsigned long>(getpid()) << '\n';

    ShmChannel channel;
    if (!channel.open(shm_name)) {
        std::cerr << "[Child] Failed to open shared memory: " << shm_name << '\n';
        return 1;
    }

    std::string line{};
    bool is_eof = false;

    while (channel.receive_line(line, is_eof)) {
        reverse_string(line);
        if (!(std::cout << line << '\n')) {
            std::cerr << "[Child] Failed to write to output file\n";
            channel.close();
            return 1;
        }
    }

    channel.close();

    if (!std::cout.flush() || !is_eof) {
        return 1;
    }

    return 0;
}
\end{lstlisting}
\begin{center}
    \small Листинг 5 — Прием строк из разделяемой памяти и их инвертирование (\path{src/child.cpp})
\end{center}

\section{Результаты работы}

\subsection{Консольный вывод и логи выполнения}
Ниже приведен фактический протокол работы программы в терминале с отображением идентификаторов процессов (\texttt{PID}), подтверждающий конкурентное выполнение родительского и двух дочерних процессов, корректную маршрутизацию и инвертирование строк по правилу длины, а также отсутствие неудаленных сегментов в \path{/dev/shm} после завершения работы:

\begin{lstlisting}
$ ./parent
[Parent] PID: 562
Enter output file for Child1: short_lines.txt
Enter output file for Child2: long_lines.txt
[Parent] Created Child1 (PID: 565)
[Parent] Created Child2 (PID: 566)
Enter lines for processing (Ctrl+D / Ctrl+C to finish): 
[Child] PID: 565
[Child] PID: 566
12345
1234567890
12345678901
123456789101112
12
hello world!
[Parent] Closed Child1
[Parent] Closed Child2

$ cat short_lines.txt
54321
0987654321
21

$ cat long_lines.txt
10987654321
211101987654321
!dlrow olleh

$ ls -la /dev/shm
total 0
drwxrwxrwt  2 root root   40 Nov 26 12:30 .
drwxr-xr-x 15 root root 3840 Nov 26 12:20 ..
\end{lstlisting}

\section{Выводы}
В результате выполнения лабораторной работы реализован программный комплекс межпроцессного взаимодействия для ОС Linux (POSIX), осуществляющий потоковую фильтрацию и инвертирование строк произвольной длины с использованием отображаемой в память разделяемой памяти (\texttt{shm\_open} + \texttt{mmap}) и неименованных POSIX-семафоров. В ходе разработки были учтены особенности модели виртуальной памяти: независимость виртуальных адресных пространств процессов, необходимость плоского представления данных без указателей, контроль границ отображаемого объекта через \texttt{fstat()} и разделение операций снятия отображения (\texttt{munmap}) и удаления имени объекта (\texttt{shm\_unlink}).

В процессе работы мною были изучены и отработаны практические навыки организации синхронизированного обмена данными через общую память без состояний гонки и взаимоблокировок, обработки сигналов прерывания и безопасного освобождения ресурсов операционной системы.

\end{document}
